% Zone „Das kennst du schon“ – aus bank/kenngroessen-von-verteilungen/zone.jsonl (zusammenbau v0.1)
\einheitenkopf[zone][Kennst du schon]{Das kennst du schon}
\setcounter{aufgabe}{0}

%% TODO Titel: Ich-kann-Satz fehlt in der Bank (Platzhalter: Kettenname) – Nr. 1
%% TODO Anweisung: ein Satz über der ersten Teilaufgabe fehlt in der Bank – Nr. 1
\begin{aufgabe}{Verteilungstabelle einer Zufallsgröße aufstellen und lesen}
\begin{teile}
\teil Die Tabelle gehört zur Zufallsgröße $X$. Lies ab und addiere: Wie groß ist $P(X \ge 2)$?

\sachtabelle{lccc}{$x$ & 1 & 2 & 3}{$P(X = x)$ & 0{,}1 & 0{,}6 & 0{,}3}
$P(X \ge 2) = $ \leerfeld
\teil Ein Glücksrad hat drei Sektoren: 72° mit der Zahl 3, 108° mit der Zahl 1 und 180° mit der Zahl 0. $X$ ist die gedrehte Zahl. Wie groß ist $P(X = 1)$? $P(X = 1) = $ \leerfeld
\end{teile}
\end{aufgabe}

%% TODO Titel: Ich-kann-Satz fehlt in der Bank (Platzhalter: Kettenname) – Nr. 2
%% TODO Anweisung: ein Satz über der ersten Teilaufgabe fehlt in der Bank – Nr. 2
\begin{aufgabe}{Pfadregeln und Gegenereignis (mit und ohne Zurücklegen)}
\begin{teile}
\teil Eine Münze wird zweimal geworfen. Wie groß ist die Wahrscheinlichkeit für zweimal Zahl? $P = $ \leerfeld
\teil In einer Urne liegen vier rote und zwei blaue Kugeln. Es werden zwei Kugeln ohne Zurücklegen gezogen. Wie groß ist die Wahrscheinlichkeit für zweimal Rot? $P = $ \leerfeld
\end{teile}
\end{aufgabe}

%% TODO Titel: Ich-kann-Satz fehlt in der Bank (Platzhalter: Kettenname) – Nr. 3
%% TODO Anweisung: ein Satz über der ersten Teilaufgabe fehlt in der Bank – Nr. 3
\begin{aufgabe}{Bernoulli-Kette erkennen}
\begin{teile}
\teil Eine Münze wird sechsmal geworfen, gezählt wird, wie oft Zahl fällt. Ist das eine Bernoulli-Kette? \\ \kreuz{ja} \\ \kreuz{nein}
\teil In einer Stadt fahren 30\,\% der Beschäftigten mit dem Rad zur Arbeit. Es werden 40 Beschäftigte zufällig befragt; Treffer ist „fährt Rad“. Wie lang ist die Kette, wie groß die Trefferwahrscheinlichkeit? $n = $ \leerfeld , $p = $ \leerfeld
\end{teile}
\end{aufgabe}

%% TODO Titel: Ich-kann-Satz fehlt in der Bank (Platzhalter: Kettenname) – Nr. 4
%% TODO Anweisung: ein Satz über der ersten Teilaufgabe fehlt in der Bank – Nr. 4
\begin{aufgabe}{Brüche gewichten und addieren, Prozentrechnung}
\begin{teile}
\teil Berechne $\frac{1}{2} \cdot 4 + \frac{1}{4} \cdot 8$.
\teil Berechne $0{,}3 \cdot 5 + 0{,}6 \cdot 2 + 0{,}1 \cdot 0$.
\end{teile}
\end{aufgabe}

%% TODO Titel: Ich-kann-Satz fehlt in der Bank (Platzhalter: Kettenname) – Nr. 5
%% TODO Anweisung: ein Satz über der ersten Teilaufgabe fehlt in der Bank – Nr. 5
\begin{aufgabe}{Lineare und quadratische Gleichungen sowie einfache Ungleichungen lösen, Lösungen sieben}
\begin{gleichungsraster}[2]
\gl{\text{Löse $3x + 2 = 14$.}} &
\gl{\text{Bestimme die kleinste natürliche Zahl $n$ mit $0{,}4 \cdot n > 7$.}} \\
\end{gleichungsraster}
\end{aufgabe}

%% TODO Titel: Ich-kann-Satz fehlt in der Bank (Platzhalter: Kettenname) – Nr. 6
%% TODO Anweisung: ein Satz über der ersten Teilaufgabe fehlt in der Bank – Nr. 6
\begin{aufgabe}{Wurzelterme bilden und vergleichen}
\begin{teile}
\teil Berechne $\sqrt{16 \cdot 4}$.
\teil Berechne $\sqrt{100 \cdot 0{,}3 \cdot 0{,}7}$ und runde auf zwei Stellen.
\end{teile}
\end{aufgabe}

%% TODO Titel: Ich-kann-Satz fehlt in der Bank (Platzhalter: Kettenname) – Nr. 7
%% TODO Anweisung: ein Satz über der ersten Teilaufgabe fehlt in der Bank – Nr. 7
\begin{aufgabe}{Pfadregeln und Gegenereignis (mit und ohne Zurücklegen) – Fehler finden}
\begin{teile}
\teil Lea soll die Wahrscheinlichkeit für genau eine rote Kugel berechnen. In der Urne liegen zwei rote und drei weiße Kugeln, gezogen wird zweimal mit Zurücklegen. Finde den Fehler und rechne richtig. \rechnung{P &= \frac{2}{5} \cdot \frac{3}{5} \\ &= \frac{6}{25}}
\end{teile}
\end{aufgabe}

%% TODO Titel: Ich-kann-Satz fehlt in der Bank (Platzhalter: Kettenname) – Nr. 8
%% TODO Anweisung: ein Satz über der ersten Teilaufgabe fehlt in der Bank – Nr. 8
\begin{aufgabe}{Pfadregeln und Gegenereignis (mit und ohne Zurücklegen) – selbst rechnen}
\begin{teile}
\teil Ein Glücksrad zeigt mit der Wahrscheinlichkeit $\frac{1}{3}$ Gelb. Es wird zweimal gedreht. Wie groß ist die Wahrscheinlichkeit für genau einmal Gelb? $P = $ \leerfeld
\end{teile}
\end{aufgabe}
